\section{Case Study: The MGM Resorts Cyberattack (September 2023)}

\subsection{Background and Context}

\textbf{Threat Actors:} Scattered Spider (UNC3944) in coordination with ALPHV/BlackCat ransomware group \\
\textbf{Attack Type:} Structured, Direct — Social Engineering combined with Ransomware-as-a-Service (RaaS) \\
\textbf{Primary Motivation:} Financial extortion; later escalated due to what the attackers described as retaliation for failed negotiations \\
\textbf{Duration:} September 8--19, 2023 — ten days of active disruption \\
\textbf{Scale:} 29 hotel and resort properties affected worldwide, including Bellagio, MGM Grand, Mandalay Bay, and Cosmopolitan in Las Vegas. Estimated \$100 million loss in third-quarter 2023 revenue. A \$45 million class-action settlement was reached in January 2025.

The MGM Resorts attack is one of the most instructive case studies in modern cybersecurity precisely because of what it did \emph{not} use. There was no zero-day vulnerability, no nation-state malware toolkit, no brute-force password cracking. All the attackers did to compromise a company valued at \$33.9 billion was search LinkedIn, find an employee, and call the help desk. Cyber Security Hub The entire entry into one of the world's largest hospitality companies cost the attackers approximately ten minutes of conversation.

This case is also notable for its two-actor architecture. Scattered Spider and ALPHV/BlackCat are distinct organizations with complementary capabilities — one specializes in social engineering to obtain access, the other provides professional-grade ransomware infrastructure. Their collaboration represents what the cybersecurity industry now calls a ``criminal supply chain,'' a division of labor that allows non-technical attackers to cause deeply technical damage.



\subsection{The Two Threat Actors}

Understanding who executed this attack is essential to understanding how it worked.

Scattered Spider (UNC3944) is a loosely organized group, mostly made up of teenagers and young adults believed to live in the United States and the United Kingdom, believed to be affiliated with the cybercriminal network known as ``The Com.'' Wikipedia Their defining capability is fluency in English and expertise in social engineering — specifically, the ability to convincingly impersonate employees over the phone to manipulate IT help desk staff. They are not primarily technical hackers; they are manipulators who exploit organizational trust structures.

ALPHV/BlackCat is a professional ransomware-as-a-service operation. BlackCat/ALPHV is a RaaS operation — just as commercial businesses rely on SaaS applications to do their work, criminal gangs rely on RaaS. BlackCat provides the professional services that Scattered Spider lacks. CyberArk The relationship between the two groups follows a division of labor: Scattered Spider gains the access, ALPHV supplies the ransomware and handles extortion operations.

\subsection{Phase 1 — Reconnaissance and Probing}

\subsubsection{The Intelligence Gathering}

The reconnaissance phase of the MGM attack required no special tools, no network scanning, and no malware. It required only a LinkedIn account and the knowledge of what to look for.

Scattered Spider used LinkedIn to identify a current MGM Resorts employee, assumed their identity, and called the MGM IT help desk requesting assistance logging into their account. SpecopsSoft From LinkedIn, the attacker gathered the employee's full name, job title, department, and likely their organizational reporting structure — all publicly available information that any recruiter or sales professional would access routinely. This is what security researchers call Open Source Intelligence (OSINT): the systematic exploitation of publicly available information to build an attack profile.

The target was selected deliberately. Scattered Spider did not pick a random employee — they identified someone whose role and credentials would grant useful access once impersonated. This represents a key insight about modern reconnaissance: the most dangerous intelligence gathering today does not happen on the dark web or through network intrusion. It happens on professional networking platforms that organizations actively encourage employees to use.

\subsubsection{Why This Phase Is the Most Critical}

Section 5.4.1 of this report identifies reconnaissance as ``the most important phase'' because everything that follows depends on its quality. The MGM attack demonstrates this in an extreme form. The entire technical complexity of the subsequent attack — the Okta takeover, the Azure Active Directory compromise, the deployment of ransomware across hundreds of servers — was made possible by ten minutes of social engineering built on a LinkedIn search.

The phone call lasted ten minutes, and the attackers were able to gain administrator privileges to MGM's Okta and Azure tenant environments. University of Hawai'i--West O'ahu No malware was involved in Phase 1. No network was probed. The reconnaissance was entirely human-focused, and its product was the single most valuable thing an attacker can obtain: legitimate administrative credentials willingly handed over by the target organization itself.

\subsection{Access and Privilege Escalation}

\subsubsection{The Entry Point: Okta}

Okta is an Identity and Access Management (IAM) platform. Organizations use it as a centralized gateway that controls authentication for all their applications and services — email, cloud infrastructure, internal tools, and more. Gaining administrator access to an organization's Okta environment is functionally equivalent to obtaining a master key to every door in a building.

This breach provided access to the Okta platform, a crucial access management system. The attackers then capitalized on weak multi-factor authentication to escalate privileges and gain control of the Azure Active Directory domain controller. Brown \& Brown The Azure Active Directory domain controller is the core identity authority for Microsoft-based cloud environments. Controlling it means controlling the ability to create, modify, and impersonate any user account in the organization.

ALPHV publicly confirmed that the affiliate group had achieved persistence in the network and had gained super administrator privileges. Morphisec With super administrator access to both the identity management layer (Okta) and the cloud directory (Azure AD), the attackers had effectively inherited the digital keys to MGM's entire infrastructure.

\subsubsection{Lateral Movement: From Cloud to Physical Operations}

What makes this attack particularly instructive is how cloud-based access translated into physical-world disruption. The IdP solution granted highly privileged access to Azure, which ultimately allowed the cloud-originated attack into MGM's brick-and-mortar operations. CyberArk The identity platform that governed digital logins was configured in a way that allowed the attackers to pivot from the cloud environment into MGM's on-premises virtualization infrastructure.

The specific mechanism involved credential synchronization between on-premise Active Directory and the cloud. Because MGM synchronized credentials between these environments, gaining control of the cloud identity layer gave Scattered Spider a pathway into the VMware ESXi hypervisor infrastructure — the virtual machine hosts that ran the servers powering MGM's hotel operations, gaming systems, and booking platforms.

The attack surface progression can be visualized as an identity compromise that propagated through cloud control into on-premises operations:

\begin{figure}[htbp]
\centering
\begin{tikzpicture}[
  node distance=1.15cm and 1.85cm,
  every node/.style={font=\sffamily\footnotesize},
  layer/.style={
    rectangle, rounded corners=4pt,
    minimum width=3.35cm, minimum height=1.62cm,
    text width=3.05cm, align=center,
    draw=black!75, line width=0.9pt
  },
  steplabel/.style={
    circle,
    minimum size=0.46cm,
    inner sep=0pt,
    draw=black!60,
    fill=white,
    line width=0.7pt,
    font=\sffamily\bfseries\tiny,
    text=black!80
  },
  flow/.style={-{Stealth[length=2.8mm,width=2.1mm]}, line width=0.95pt, draw=black!75},
  flowlabel/.style={
    font=\sffamily\scriptsize\bfseries,
    fill=white,
    rounded corners=1.5pt,
    inner xsep=2.1pt,
    inner ysep=1.2pt,
    text=black!85
  }
]

\node[layer, fill=violet!24] (identity) {\textbf{Identity layer}\\Okta + Azure AD\\Admin access seized};
\node[layer, fill=orange!35, right=of identity] (cloud) {\textbf{Cloud control}\\SSO + tenant admin\\Privilege expansion};
\node[layer, fill=teal!30, right=of cloud] (onprem) {\textbf{On-prem operations}\\VMware ESXi + hotel systems\\Ransomware impact zone};

\node[steplabel, above=3pt of identity] {1};
\node[steplabel, above=3pt of cloud] {2};
\node[steplabel, above=3pt of onprem] {3};

\draw[flow] (identity) -- node[midway, above=3pt, flowlabel] {Credential abuse} (cloud);
\draw[flow] (cloud) -- node[midway, below=3pt, flowlabel] {Sync + pivot path} (onprem);

\end{tikzpicture}
\caption{Attack surface progression from identity compromise to on-prem operational impact}
\label{fig:mgm-attack-surface-map}
\end{figure}

\subsubsection{Data Exfiltration: The Double-Extortion Setup}

Before deploying any ransomware, the attackers ensured they had a second lever of pressure. ALPHV claims to have exfiltrated 6 TB of customer information during this time, upon which they initiated negotiations with MGM to prevent the public release of the stolen data. ALPHV also threatened to disclose the exfiltrated information if an agreement could not be reached. Netwrix

This sequencing is deliberate and important. Data theft before encryption means that even if the victim restores all systems from backup — eliminating the encryption leverage entirely — the attacker retains the ability to cause reputational and regulatory damage through data publication. This dual-pressure structure is what the industry now calls ``double extortion,'' and it directly undermines the traditional defense of maintaining strong backups.

\subsubsection{The Ransomware Deployment}

At this stage, ALPHV/BlackCat took operational control. Using the BlackCat/ALPHV RaaS service, Scattered Spider encrypted several hundred ESXi servers, which hosted thousands of virtual machines supporting hundreds of systems widely used in the hospitality industry. This caused cascading chaos. As the ESXi hosts became encrypted one after another, the applications running on them crashed. CyberArk

The operational consequences were immediate and visible to the public. The effects of the cyberattack on MGM Resorts were immediately apparent: gambling operations were disrupted as slot machines went offline, hotel guests found their digital room keys had stopped working, online reservation and booking systems were shut down, and the MGM mobile app became completely inaccessible. Netwrix Staff reverted to pen and paper to manage guest queues. The attack had successfully converted a digital breach into a physical operational crisis.

\subsection{Phase 3 — Covering Traces of the Attack}

\subsubsection{The Backdoor Architecture}

The most forensically significant aspect of Phase 3 is that it began before MGM detected anything. The ransomware was deployed after MGM locked out the network, which means that the attackers had gained significant visibility of the network and had implanted backdoors across the network before the main attack was triggered. Morphisec

This is a critical point that separates sophisticated attackers from unsophisticated ones. By planting backdoors before becoming visible, the attackers ensured that MGM's detection and response would not equal eviction. Even if MGM successfully expelled the attackers from their primary access route, the pre-planted backdoors would provide re-entry paths. This is the attack equivalent of having multiple escape routes prepared before a heist begins.

\subsubsection{Escalation Upon Detection: Attacking the Responders}

MGM's incident response team identified the intrusion the day after initial access was gained — September 9 — when they detected unusual activity and traffic on the company's systems. Their response was to terminate the Okta Sync servers, which was the most direct way to cut off the attacker's credential-harvesting operations. This response was tactically correct but strategically insufficient.

ALPHV admitted to sniffing passwords on their Okta servers after the initial access. MGM hastily deactivated their Okta Sync servers and essential infrastructure components to prevent an escalation of the attack, causing the interruption of reservation systems, digital room keys, slot machines, and more. University of Hawai'i--West O'ahu The shutting down of Okta was itself what caused the guest-facing outages — meaning MGM's defensive action created the public disruption, not the ransomware directly.

In a statement, BlackCat noted that MGM Resorts had not responded through the provided communication channel, indicating a reluctance to negotiate. The hackers highlighted that MGM's only response to the breach had been to disconnect each and every one of their Okta Sync servers after discovering the intrusion. Avertium The attackers, still possessing access through their pre-planted backdoors, then deployed the ransomware as an escalation — consistent with their own statement that the ransomware component was partly driven by retaliation for what they perceived as bad-faith negotiation.

\subsubsection{Attribution and Legal Aftermath}

Despite the attackers' technical evasion during the operation, the group's English-speaking, socially-visible profile made eventual attribution feasible. In July 2024, the West Midlands Police with the help of the FBI arrested a 17-year-old juvenile in connection with the MGM cyberattacks. In November 2024, 19-year-old Remington Ogletree was arrested on charges related to his alleged involvement with the group. Wikipedia

The arrests underscore a pattern: sophisticated operational security at the technical level does not compensate for poor personal operational security. Members of Scattered Spider communicated openly on Telegram channels and were identifiable through social media activity, which ultimately provided law enforcement with the leads needed to locate them despite the digital obfuscation during the attack itself.'


Here is the complete attack chain from reconnaissance through ransomware deployment:

\begin{figure}[htbp]
\centering
\begin{tikzpicture}[
  node distance = 0.9cm,
  every node/.style = {font=\sffamily\small},
  startstop/.style = {
    rectangle, rounded corners=3pt,
    minimum width=8.4cm, minimum height=0.85cm,
    text centered, draw=black!75, fill=gray!18,
    line width=0.9pt, align=center, font=\bfseries
  },
  process/.style = {
    rectangle, rounded corners=3pt,
    minimum width=8.4cm, minimum height=1.08cm,
    text width=8.0cm, text centered, draw=black!75,
    line width=0.8pt, align=center
  },
  phase1/.style = {process, fill=violet!24},
  phase2/.style = {process, fill=orange!42},
  phase3/.style = {process, fill=teal!34},
  arrow/.style = {-{Stealth[length=2.7mm,width=2mm]}, line width=0.9pt, draw=black!75},
  legendbox/.style = {
    rectangle, rounded corners=3pt,
    minimum width=3.3cm, minimum height=0.9cm,
    draw=black!45, line width=0.7pt,
    align=center, font=\sffamily\bfseries\footnotesize
  },
  actorlabel/.style = {font=\bfseries\footnotesize, anchor=west}
]

% Legend (balanced and centered above the timeline)
\node[legendbox, fill=violet!24] (leg1) {Phase 1\\Reconnaissance};
\node[legendbox, fill=orange!42, right=0.25cm of leg1] (leg2) {Phase 2\\Access \& escalation};
\node[legendbox, fill=teal!34, right=0.25cm of leg2] (leg3) {Phase 3\\Covering traces};

% Main vertical chain
\node[startstop, below=0.95cm of leg2] (start) {Sep 7, 2023 --- attack begins};

\node[phase1, below=of start] (osint) {\textbf{LinkedIn OSINT} \\ Target employee identified by name + role};

\node[phase1, below=of osint] (vishing) {\textbf{IT help desk vishing} \\ 10-min call — admin credentials handed over};

\node[phase2, below=of vishing] (okta) {\textbf{Okta + Azure AD takeover} \\ Super admin privileges gained};

\node[phase2, below=of okta] (lateral) {\textbf{Lateral movement} \\ 6 TB customer data exfiltrated};

\node[phase3, below=of lateral] (backdoors) {\textbf{Backdoors planted} \\ MGM detects — shuts down Okta};

\node[phase3, below=of backdoors] (ransom) {\textbf{ALPHV ransomware deployed} \\ 100+ ESXi servers encrypted};

\node[startstop, below=of ransom] (end) {Sep 20 — Systems restored};

% Arrows between nodes
\draw[arrow] (start) -- (osint);
\draw[arrow] (osint) -- (vishing);
\draw[arrow] (vishing) -- (okta);
\draw[arrow] (okta) -- (lateral);
\draw[arrow] (lateral) -- (backdoors);
\draw[arrow] (backdoors) -- (ransom);
\draw[arrow] (ransom) -- (end);

% Right-side actor grouping labels + subtle grouping lines
% Use overlay so these annotations do not affect centering/bounding box
\coordinate (actorX) at ($(lateral.east)+(1.45,0)$);

\draw[overlay, gray!55, line width=0.8pt] (actorX |- osint.north) -- (actorX |- lateral.south);
\node[overlay, actorlabel] at ($(actorX)+(0.18,0.95)$) {Scattered Spider};

\draw[overlay, gray!55, line width=0.8pt] (actorX |- backdoors.north) -- (actorX |- ransom.south);
\node[overlay, actorlabel] at ($(actorX)+(0.18,-1.65)$) {ALPHV / BlackCat};

\end{tikzpicture}
\caption{Complete attack chain of the MGM Resorts cyberattack (September 2023)}
\label{fig:mgm-attack-chain}
\end{figure}

\begin{figure}[htbp]
\centering
\begin{tikzpicture}[
  x=0.86cm,
  y=1cm,
  every node/.style={font=\sffamily\scriptsize},
  milestone/.style={circle, fill=black!75, inner sep=1.6pt},
  calloutTop/.style={
    rectangle, rounded corners=2pt,
    draw=black!60, fill=white,
    align=center, inner sep=3pt,
    text width=2.45cm,
    anchor=south
  },
  calloutBottom/.style={
    rectangle, rounded corners=2pt,
    draw=black!60, fill=white,
    align=center, inner sep=3pt,
    text width=2.45cm,
    anchor=north
  },
  phasebar/.style={
    rounded corners=2pt,
    line width=0.7pt,
    align=center,
    font=\sffamily\scriptsize
  }
]

% Base timeline
\draw[thick, -{Stealth[length=2.5mm,width=2mm]}] (0,0) -- (11.45,0);
\node[font=\sffamily\bfseries\footnotesize] at (5.6,2.0) {September 2023 response timeline};

% Day ticks (Sep 8 to Sep 19)
\foreach \x/\d in {0/8,1/9,2/10,3/11,4/12,5/13,6/14,7/15,8/16,9/17,10/18,11/19} {
  \draw (\x,0.1) -- (\x,-0.1);
  \node[below=3pt] at (\x,-0.1) {\d};
}

% Milestones and events
\node[milestone] (d8) at (0,0) {};
\node[calloutTop] (c8) at (0,0.98) {Sep 8\\Initial access\\(vishing success)};
\draw[gray!70, -{Stealth[length=2mm,width=1.6mm]}] (c8.south) -- (d8.north);

\node[milestone] (d9) at (1,0) {};
\node[calloutBottom] (c9) at (1,-0.98) {Sep 9\\Detection +\\Okta Sync off};
\draw[gray!70, -{Stealth[length=2mm,width=1.6mm]}] (c9.north) -- (d9.south);

\node[phasebar, fill=orange!28, draw=orange!70, anchor=south, minimum height=0.44cm, text width=2.9cm]
  at (3.5,0.43) {Sep 10--12\\Lateral movement\\+ 6 TB exfiltration};

\node[milestone] (d13) at (5,0) {};
\node[calloutTop] (c13) at (5,1.08) {Sep 13\\ALPHV ransomware\\deployed};
\draw[gray!70, -{Stealth[length=2mm,width=1.6mm]}] (c13.south) -- (d13.north);

\node[phasebar, fill=teal!24, draw=teal!70, anchor=north, minimum height=0.44cm, text width=3.25cm]
  at (8.45,-0.52) {Sep 14--18\\Containment +\\staged recovery};

\node[milestone] (d19) at (11,0) {};
\node[calloutTop] (c19) at (11,0.98) {Sep 19\\Core systems\\restored};
\draw[gray!70, -{Stealth[length=2mm,width=1.6mm]}] (c19.south) -- (d19.north);

\end{tikzpicture}
\caption{Day-by-day response timeline of the MGM incident (Sep 8--19, 2023)}
\label{fig:mgm-response-timeline}
\end{figure}



\subsection{Impact Analysis}

The operational consequences of the attack were severe and immediate across every dimension of MGM's business.

Financial losses totaled approximately \$100 million in third-quarter 2023 results, with MGM reporting daily revenue losses of roughly \$8.4 million during the ten-day outage. The company additionally committed \$40 million in new cybersecurity investments following the incident, and agreed to a \$45 million class-action settlement in January 2025 covering affected customers from both the 2019 and 2023 breaches.

\begin{table}[htbp]
\centering
\small
\setlength{\tabcolsep}{5pt}
\begin{tabularx}{\textwidth}{>{\bfseries}p{2.8cm} X p{3.05cm} p{2.45cm}}
\toprule
Impact area & Evidence & Estimated loss & Duration \\
\midrule
Financial impact & Q3 2023 revenue reduction and post-incident security spending commitments & \$100M reported outage impact; \$40M additional cybersecurity investment; \$45M class-action settlement & Approx. 10 days of active disruption; legal costs extending into 2025 \\

Operational impact & Slot machines, room keys, booking systems, MGM app, and guest workflows disrupted; staff shifted to manual processing & Daily loss estimated at \$8.4M during outage window & Service degradation across Sep 8--19 \\

Data exposure & Attackers claimed exfiltration of 6 TB of customer data and used leak threats in extortion & No final public total; significant legal and reputational liability exposure & Data theft occurred before encryption phase \\

Regulatory / legal & SEC incident disclosure obligations and class-action litigation outcomes & \$45M settlement (announced Jan 2025) & Post-incident consequences extending beyond restoration \\
\bottomrule
\end{tabularx}
\caption{Incident impact summary for the MGM Resorts cyberattack}
\label{tab:mgm-incident-impact-summary}
\end{table}

Operational disruption was total and publicly visible in a way that most cyberattacks are not. Slot machines displayed error messages. Digital room keys failed. Guests could not check in online, make card payments, or access their MGM loyalty accounts. Staff processed hotel transactions by hand. Airport lounges and concierge services lost access to reservation systems. For a hospitality company whose entire value proposition is seamless guest experience, this was the maximum possible reputational damage.

Regulatory and legal consequences extended beyond the settlement. MGM and Caesars Entertainment — which suffered a parallel attack by Scattered Spider the same week and paid a \$15 million ransom — were both required to file mandatory disclosures to the Securities and Exchange Commission under new cybersecurity incident rules. Moody's Corporation noted the attack may negatively affect MGM's credit rating due to the company's heavy operational reliance on technology systems.



\subsection{Mapping to the CIA Triad}

This attack methodically violated all three pillars of the CIA Triad introduced in Chapter 1 of this report.

Confidentiality was breached through the exfiltration of 6 TB of customer data, including personally identifiable information belonging to guests, employees, and vendors. The threat of public data release formed the second half of ALPHV's double-extortion strategy, maintaining confidentiality pressure even after systems were restored.

Integrity was compromised through the unauthorized modification of MGM's identity infrastructure. By gaining administrator control of Okta and Azure Active Directory, the attackers could create, alter, and impersonate accounts — corrupting the integrity of the authentication system that the entire organization depended upon.

Availability was the most visibly destroyed pillar. The encryption of over one hundred ESXi hypervisors hosting thousands of virtual machines rendered MGM's core operational systems — reservations, gaming, payments, room access — completely unavailable for ten days, directly harming guests and generating tens of millions of dollars in lost revenue.

\begin{table}[htbp]
\centering
\small
\setlength{\tabcolsep}{5pt}
\begin{tabularx}{\textwidth}{>{\bfseries}p{2.8cm} X X}
\toprule
CIA pillar & What was affected & Concrete MGM example \\
\midrule
Confidentiality & Sensitive customer and organizational data secrecy & Claimed exfiltration of 6 TB of customer information, followed by threat of public release in extortion pressure \\

Integrity & Trustworthiness of identity and account-control systems & Unauthorized admin control over Okta and Azure AD enabled account manipulation and impersonation capability \\

Availability & Continuous access to critical business services & Encryption of 100+ ESXi hosts disrupted reservations, gaming, payments, room access, and mobile services for roughly ten days \\
\bottomrule
\end{tabularx}
\caption{CIA triad mapping for the MGM Resorts incident}
\label{tab:mgm-cia-triad-mapping}
\end{table}



\subsection{Defensive Lessons}

The MGM attack exposes failures at multiple layers of what Section 7 of this report describes as Defense in Depth. Each layer that should have stopped or slowed the attack either did not exist or was insufficient.

The help desk had no caller verification protocol. The single most important lesson from this attack is that a ten-minute phone call should never produce administrator credentials. Secure service desk practices require callers to authenticate with something they have — a registered device, a hardware token, a one-time code — not only something they know, such as their name and job title. An attacker with LinkedIn access knows everything that knowledge-based verification checks for.

Multi-factor authentication on privileged accounts was weak. The attackers capitalized on weak multi-factor authentication to escalate privileges and gain control of the Azure Active Directory domain controller. Brown \& Brown MFA on administrative accounts must be phishing-resistant — hardware security keys or certificate-based authentication — not SMS or app-push approvals, both of which are vulnerable to social engineering.

The network lacked segmentation between identity and operations. The critical failure that converted a digital breach into a physical operational crisis was the pathway from cloud identity management to on-premise ESXi infrastructure. A properly segmented network with strict boundaries between the identity management layer, corporate IT, and operational technology would have contained the blast radius significantly.

Backups did not protect against the double-extortion model. MGM's ability to restore systems after ten days demonstrates that backup capabilities existed. However, restoring data from backup does nothing to address the 6 TB of already-exfiltrated customer records. The triple-extortion model described in Chapter 8 of this report — encryption, data theft, and harassment — renders traditional backup-as-defense incomplete.

User awareness training must explicitly address vishing. The IT help desk staff who granted credentials were not acting maliciously or carelessly by their own training standards — they were following normal helpdesk procedures. User awareness programs must include specific scenarios for voice-based impersonation attacks, with clear protocols that override normal helpfulness instincts when administrative access is at stake.